\documentclass[a4paper,12pt,twoside]{memoir}

% Castellano
\usepackage[spanish,es-tabla]{babel}
\selectlanguage{spanish}
\usepackage[utf8]{inputenc}
\usepackage[T1]{fontenc}
\usepackage{lmodern} % Scalable font
\usepackage{microtype}
\usepackage{placeins}

\RequirePackage{booktabs}
\RequirePackage[table]{xcolor}
\RequirePackage{xtab}
\RequirePackage{multirow}

% Links
\PassOptionsToPackage{hyphens}{url}\usepackage[colorlinks]{hyperref}
\hypersetup{
	allcolors = {red}
}

% Ecuaciones
\usepackage{amsmath}

% Rutas de fichero / paquete
\newcommand{\ruta}[1]{{\sffamily #1}}

% Párrafos
\nonzeroparskip

% Huérfanas y viudas
\widowpenalty100000
\clubpenalty100000

% Imágenes

% Comando para insertar una imagen en un lugar concreto.
% Los parámetros son:
% 1 --> Ruta absoluta/relativa de la figura
% 2 --> Texto a pie de figura
% 3 --> Tamaño en tanto por uno relativo al ancho de página
\usepackage{graphicx}
\newcommand{\imagen}[3]{
	\begin{figure}[!h]
		\centering
		\includegraphics[width=#3\textwidth]{#1}
		\caption{#2}\label{fig:#1}
	\end{figure}
	\FloatBarrier
}

% Comando para insertar una imagen sin posición.
% Los parámetros son:
% 1 --> Ruta absoluta/relativa de la figura
% 2 --> Texto a pie de figura
% 3 --> Tamaño en tanto por uno relativo al ancho de página
\newcommand{\imagenflotante}[3]{
	\begin{figure}
		\centering
		\includegraphics[width=#3\textwidth]{#1}
		\caption{#2}\label{fig:#1}
	\end{figure}
}

% El comando \figura nos permite insertar figuras comodamente, y utilizando
% siempre el mismo formato. Los parametros son:
% 1 --> Porcentaje del ancho de página que ocupará la figura (de 0 a 1)
% 2 --> Fichero de la imagen
% 3 --> Texto a pie de imagen
% 4 --> Etiqueta (label) para referencias
% 5 --> Opciones que queramos pasarle al \includegraphics
% 6 --> Opciones de posicionamiento a pasarle a \begin{figure}
\newcommand{\figuraConPosicion}[6]{%
  \setlength{\anchoFloat}{#1\textwidth}%
  \addtolength{\anchoFloat}{-4\fboxsep}%
  \setlength{\anchoFigura}{\anchoFloat}%
  \begin{figure}[#6]
    \begin{center}%
      \Ovalbox{%
        \begin{minipage}{\anchoFloat}%
          \begin{center}%
            \includegraphics[width=\anchoFigura,#5]{#2}%
            \caption{#3}%
            \label{#4}%
          \end{center}%
        \end{minipage}
      }%
    \end{center}%
  \end{figure}%
}

%
% Comando para incluir imágenes en formato apaisado (sin marco).
\newcommand{\figuraApaisadaSinMarco}[5]{%
  \begin{figure}%
    \begin{center}%
    \includegraphics[angle=90,height=#1\textheight,#5]{#2}%
    \caption{#3}%
    \label{#4}%
    \end{center}%
  \end{figure}%
}
% Para las tablas
\newcommand{\otoprule}{\midrule [\heavyrulewidth]}
%
% Nuevo comando para tablas pequeñas (menos de una página).
\newcommand{\tablaSmall}[5]{%
 \begin{table}
  \begin{center}
   \rowcolors {2}{gray!35}{}
   \begin{tabular}{#2}
    \toprule
    #4
    \otoprule
    #5
    \bottomrule
   \end{tabular}
   \caption{#1}
   \label{tabla:#3}
  \end{center}
 \end{table}
}

%
% Nuevo comando para tablas pequeñas (menos de una página).
\newcommand{\tablaSmallSinColores}[5]{%
 \begin{table}[H]
  \begin{center}
   \begin{tabular}{#2}
    \toprule
    #4
    \otoprule
    #5
    \bottomrule
   \end{tabular}
   \caption{#1}
   \label{tabla:#3}
  \end{center}
 \end{table}
}

\newcommand{\tablaApaisadaSmall}[5]{%
\begin{landscape}
  \begin{table}
   \begin{center}
    \rowcolors {2}{gray!35}{}
    \begin{tabular}{#2}
     \toprule
     #4
     \otoprule
     #5
     \bottomrule
    \end{tabular}
    \caption{#1}
    \label{tabla:#3}
   \end{center}
  \end{table}
\end{landscape}
}

%
% Nuevo comando para tablas grandes con cabecera y filas alternas coloreadas en gris.
\newcommand{\tabla}[6]{%
  \begin{center}
    \tablefirsthead{
      \toprule
      #5
      \otoprule
    }
    \tablehead{
      \multicolumn{#3}{l}{\small\sl continúa desde la página anterior}\\
      \toprule
      #5
      \otoprule
    }
    \tabletail{
      \hline
      \multicolumn{#3}{r}{\small\sl continúa en la página siguiente}\\
    }
    \tablelasttail{
      \hline
    }
    \bottomcaption{#1}
    \rowcolors {2}{gray!35}{}
    \begin{xtabular}{#2}
      #6
      \bottomrule
    \end{xtabular}
    \label{tabla:#4}
  \end{center}
}

%
% Nuevo comando para tablas grandes con cabecera.
\newcommand{\tablaSinColores}[6]{%
  \begin{center}
    \tablefirsthead{
      \toprule
      #5
      \otoprule
    }
    \tablehead{
      \multicolumn{#3}{l}{\small\sl continúa desde la página anterior}\\
      \toprule
      #5
      \otoprule
    }
    \tabletail{
      \hline
      \multicolumn{#3}{r}{\small\sl continúa en la página siguiente}\\
    }
    \tablelasttail{
      \hline
    }
    \bottomcaption{#1}
    \begin{xtabular}{#2}
      #6
      \bottomrule
    \end{xtabular}
    \label{tabla:#4}
  \end{center}
}

%
% Nuevo comando para tablas grandes sin cabecera.
\newcommand{\tablaSinCabecera}[5]{%
  \begin{center}
    \tablefirsthead{
      \toprule
    }
    \tablehead{
      \multicolumn{#3}{l}{\small\sl continúa desde la página anterior}\\
      \hline
    }
    \tabletail{
      \hline
      \multicolumn{#3}{r}{\small\sl continúa en la página siguiente}\\
    }
    \tablelasttail{
      \hline
    }
    \bottomcaption{#1}
  \begin{xtabular}{#2}
    #5
   \bottomrule
  \end{xtabular}
  \label{tabla:#4}
  \end{center}
}



\definecolor{cgoLight}{HTML}{EEEEEE}
\definecolor{cgoExtralight}{HTML}{FFFFFF}

%
% Nuevo comando para tablas grandes sin cabecera.
\newcommand{\tablaSinCabeceraConBandas}[5]{%
  \begin{center}
    \tablefirsthead{
      \toprule
    }
    \tablehead{
      \multicolumn{#3}{l}{\small\sl continúa desde la página anterior}\\
      \hline
    }
    \tabletail{
      \hline
      \multicolumn{#3}{r}{\small\sl continúa en la página siguiente}\\
    }
    \tablelasttail{
      \hline
    }
    \bottomcaption{#1}
    \rowcolors[]{1}{cgoExtralight}{cgoLight}

  \begin{xtabular}{#2}
    #5
   \bottomrule
  \end{xtabular}
  \label{tabla:#4}
  \end{center}
}



\graphicspath{ {./img/} }

% Capítulos
\chapterstyle{bianchi}
\newcommand{\capitulo}[2]{
	\setcounter{chapter}{#1}
	\setcounter{section}{0}
	\setcounter{figure}{0}
	\setcounter{table}{0}
	\chapter*{#2}
	\addcontentsline{toc}{chapter}{#2}
	\markboth{#2}{#2}
}

% Apéndices
\renewcommand{\appendixname}{Apéndice}
\renewcommand*\cftappendixname{\appendixname}

\newcommand{\apendice}[1]{
	%\renewcommand{\thechapter}{A}
	\chapter{#1}
}

\renewcommand*\cftappendixname{\appendixname\ }

% Formato de portada
\makeatletter
\usepackage{xcolor}
\newcommand{\tutor}[1]{\def\@tutor{#1}}
\newcommand{\course}[1]{\def\@course{#1}}
\definecolor{cpardoBox}{HTML}{E6E6FF}
\def\maketitle{
  \null
  \thispagestyle{empty}
  % Cabecera ----------------
\noindent\includegraphics[width=\textwidth]{cabecera}\vspace{1cm}%
  \vfill
  % Título proyecto y escudo informática ----------------
  \colorbox{cpardoBox}{%
    \begin{minipage}{.8\textwidth}
      \vspace{.5cm}\Large
      \begin{center}
      \textbf{TFG del Grado en Ingeniería Informática}\vspace{.6cm}\\
      \textbf{\LARGE\@title{}}
      \end{center}
      \vspace{.2cm}
    \end{minipage}

  }%
  \hfill\begin{minipage}{.20\textwidth}
    \includegraphics[width=\textwidth]{escudoInfor}
  \end{minipage}
  \vfill
  % Datos de alumno, curso y tutores ------------------
  \begin{center}%
  {%
    \noindent\LARGE
    Presentado por \@author{}\\ 
    en Universidad de Burgos --- \@date{}\\
    Tutor: \@tutor{}\\
  }%
  \end{center}%
  \null
  \cleardoublepage
  }
\makeatother

\newcommand{\nombre}{Nombre del alumno} %%% cambio de comando

% Datos de portada
\title{título del TFG}
\author{\nombre}
\tutor{nombre tutor}
\date{\today}

\begin{document}

\maketitle


\newpage\null\thispagestyle{empty}\newpage


%%%%%%%%%%%%%%%%%%%%%%%%%%%%%%%%%%%%%%%%%%%%%%%%%%%%%%%%%%%%%%%%%%%%%%%%%%%%%%%%%%%%%%%%
\thispagestyle{empty}


\noindent\includegraphics[width=\textwidth]{cabecera}\vspace{1cm}

\noindent D. nombre tutor, profesor del departamento de nombre departamento, área de nombre área.

\noindent Expone:

\noindent Que el alumno D. \nombre, con DNI dni, ha realizado el Trabajo final de Grado en Ingeniería Informática titulado título de TFG. 

\noindent Y que dicho trabajo ha sido realizado por el alumno bajo la dirección del que suscribe, en virtud de lo cual se autoriza su presentación y defensa.

\begin{center} %\large
En Burgos, {\large \today}
\end{center}

\vfill\vfill\vfill

% Author and supervisor
\begin{minipage}{0.45\textwidth}
\begin{flushleft} %\large
Vº. Bº. del Tutor:\\[2cm]
D. nombre tutor
\end{flushleft}
\end{minipage}
\hfill
\begin{minipage}{0.45\textwidth}
\begin{flushleft} %\large
Vº. Bº. del co-tutor:\\[2cm]
D. nombre co-tutor
\end{flushleft}
\end{minipage}
\hfill

\vfill

% para casos con solo un tutor comentar lo anterior
% y descomentar lo siguiente
%Vº. Bº. del Tutor:\\[2cm]
%D. nombre tutor


\newpage\null\thispagestyle{empty}\newpage




\frontmatter

% Abstract en castellano
\renewcommand*\abstractname{Resumen}
\begin{abstract}
En este primer apartado se hace una \textbf{breve} presentación del tema que se aborda en el proyecto.
\end{abstract}

\renewcommand*\abstractname{Descriptores}
\begin{abstract}
Palabras separadas por comas que identifiquen el contenido del proyecto Ej: servidor web, buscador de vuelos, android \ldots
\end{abstract}

\clearpage

% Abstract en inglés
\renewcommand*\abstractname{Abstract}
\begin{abstract}
A \textbf{brief} presentation of the topic addressed in the project.
\end{abstract}

\renewcommand*\abstractname{Keywords}
\begin{abstract}
keywords separated by commas.
\end{abstract}

\clearpage

% Indices
\tableofcontents

\clearpage

\listoffigures

\clearpage

\listoftables
\clearpage

\mainmatter
\capitulo{1}{Introducción}

El factor de impacto de una publicación científica mide la frecuencia con la cual ha sido citado el artículo promedio de una revista en un año en particular. Específicamente, sirve para evaluar la importancia de una revista dentro de un determinado campo científico. Existen múltiples metodologías de cálculo y sus correspondientes métricas, siendo el JCR (Journal Citation Reports) y el SJR (Scimago Journal Rank) los dos índices de impacto más utilizados.
Además, cabe destacar que el factor de impacto es uno de los principales criterios empleados en los procesos de acreditación y promoción interna para evaluar la calidad del trabajo científico de millones de académicos en todo el mundo.

Por su propia naturaleza, el factor de impacto se calcula con carácter retrospectivo, i.e., sobre datos de años anteriores. Así pues, si a la hora de seleccionar la revista a la que mandarán un nuevo trabajo, los académicos quieren tener en cuenta el posible factor de impacto que tendrá la revista en el año de publicación del artículo, lo único que pueden hacer es fijarse en su índice de impacto de los años anteriores y hacer sus propias hipótesis/predicciones de futuro.

Dado que a día de hoy la herramienta Google Scholar recoge información extremadamente actualizada sobre la publicación y citación de artículos científicos (podría decirse que se actualiza prácticamente en tiempo real), creemos que puede ser muy útil para la comunidad científica utilizar los datos de Google Scholar como \textit{inputs} de modelos de aprendizaje automático para estimar el índice de impacto que tendrán las distintas revistas científicas en el año en curso.

El \textit{output} esperado del proyecto será una aplicación web de tipo open-access, la cual implementará algoritmos de aprendizaje supervisado que utilizarán los datos históricos disponibles en Google Scholar para predecir el valor del índice de impacto de todas las revistas científicas indexadas en el JCR. Dicha aplicación se dejará en un repositorio público, para así garantizar que pueda ser utilizada por toda la comunidad científica. 

\capitulo{2}{Objetivos del proyecto}

El objetivo principal de este proyecto consiste en llevar a cabo una investigación que permita sobre cómo estimar el Índice de Impacto de una revista. Para lograrlo, se desarrollará un modelo de inteligencia artificial capaz de realizar predicciones precisas al respecto. Como último paso, se diseñará y programará una aplicación web que facilite el acceso a los resultados obtenidos.

A continuación, se enumeran los principales objetivos de este proyecto:

\begin{enumerate}
\item Lectura de literatura científica sobre bibliometría para comprender bien el marco conceptual en el que se encuadra este proyecto.
\item Estudio de la metodología de cálculo de los distintos índices de impacto en general, y del JCR (Journal Citation Report) en particular. Asimismo, se estudiarán las sucesivas modificaciones/excepciones que se han ido introduciendo en el cálculo del JCR a lo largo de los años.
\item Estudio de la API de Google Scholar y otras APIs.
\item Diseño y creación de una base de datos en la que se almacenará la información bibliográfica descargada.
\item Implementación de las funciones de cálculo del índice JCR para aplicarlas sobre los datos extraídos. 
\item Experimentación y evaluación de distintos modelos de regresión (aprendizaje supervisado) para predecir el índice de impacto JCR a partir de las series temporales históricas disponibles. Selección del mejor algoritmo.
\item Diseño y creación de una aplicación web donde se incluirá el mejor modelo de regresión para ponerlo a disposición de la comunidad científica.
\end{enumerate}
\capitulo{3}{Conceptos teóricos}

En aquellos proyectos que necesiten para su comprensión y desarrollo de unos conceptos teóricos de una determinada materia o de un determinado dominio de conocimiento, debe existir un apartado que sintetice dichos conceptos.

Algunos conceptos teóricos de \LaTeX \footnote{Créditos a los proyectos de Álvaro López Cantero: Configurador de Presupuestos y Roberto Izquierdo Amo: PLQuiz}.

\section{Secciones}

Las secciones se incluyen con el comando section.

\subsection{Subsecciones}

Además de secciones tenemos subsecciones.

\subsubsection{Subsubsecciones}

Y subsecciones. 


\section{Referencias}

Las referencias se incluyen en el texto usando cite \cite{wiki:latex}. Para citar webs, artículos o libros \cite{koza92}.


\section{Imágenes}

Se pueden incluir imágenes con los comandos standard de \LaTeX, pero esta plantilla dispone de comandos propios como por ejemplo el siguiente:

\imagen{escudoInfor}{Autómata para una expresión vacía}



\section{Listas de items}

Existen tres posibilidades:

\begin{itemize}
	\item primer item.
	\item segundo item.
\end{itemize}

\begin{enumerate}
	\item primer item.
	\item segundo item.
\end{enumerate}

\begin{description}
	\item[Primer item] más información sobre el primer item.
	\item[Segundo item] más información sobre el segundo item.
\end{description}
	
\begin{itemize}
\item 
\end{itemize}

\section{Tablas}

Igualmente se pueden usar los comandos específicos de \LaTeX o bien usar alguno de los comandos de la plantilla.

\tablaSmall{Herramientas y tecnologías utilizadas en cada parte del proyecto}{l c c c c}{herramientasportipodeuso}
{ \multicolumn{1}{l}{Herramientas} & App AngularJS & API REST & BD & Memoria \\}{ 
HTML5 & X & & &\\
CSS3 & X & & &\\
BOOTSTRAP & X & & &\\
JavaScript & X & & &\\
AngularJS & X & & &\\
Bower & X & & &\\
PHP & & X & &\\
Karma + Jasmine & X & & &\\
Slim framework & & X & &\\
Idiorm & & X & &\\
Composer & & X & &\\
JSON & X & X & &\\
PhpStorm & X & X & &\\
MySQL & & & X &\\
PhpMyAdmin & & & X &\\
Git + BitBucket & X & X & X & X\\
Mik\TeX{} & & & & X\\
\TeX{}Maker & & & & X\\
Astah & & & & X\\
Balsamiq Mockups & X & & &\\
VersionOne & X & X & X & X\\
} 

\capitulo{4}{Técnicas y herramientas}

En esta sección se presentarán las distintas herramientas y recursos que se han utilizado para la realización del proyecto. 
% aquí debería venir una pequeña enumeración de lenguajes, bibliotecas, APIs, BBDD...


\section{Técnicas}
En el presente proyecto, se han utilizado diversas técnicas para llevar a cabo el análisis y cálculo del índice de impacto de publicaciones científicas. Estas técnicas han sido seleccionadas con el objetivo de obtener resultados precisos y confiables, y de cubrir las necesidades específicas de este estudio.

\subsection{Web-Scraping}
Es la práctica de recopilar datos a través de un programa que interactúe con una API~\cite{mitchell2018}. Más concretamente, un programa automatizado compuesto por \textit{queries} que realizan solicitudes HTTP para adquirir recursos de un sitio web específico. Esta solicitud se puede formatear en una URL que contenga una consulta GET o en un mensaje HTTP que contenga una consulta POST~\cite{zhao2017}. Una vez que la petición es exitosamente recibida y manejada por el sitio web seleccionado, el recurso requerido será extraído y luego devuelto al programa de \textit{web-scraping} específico.

El uso \textit{web-scraping} resulta ser una técnica efectiva en este proyecto debido a que permite obtener datos de las distintas fuentes web de manera automatizada y eficiente. La información obtenida a través de esta técnica es esencial para calcular el índice de impacto (JCR). Además, el \textit{web-scraping} permite obtener grandes cantidades de información en un corto período de tiempo, lo que resulta muy útil en proyectos en el que se requiere una gran cantidad de datos.


\section{Lenguaje de programación}
En primer lugar, nos planteamos la cuestión del lenguaje o lenguajes de programación más adecuados para nuestro objetivo. Las listas de popularidad actuales nos muestran dos entornos ganadores para proyectos de \textit{web-scraping}: Python y Javascript.

Aunque ambos lenguajes son altamente capaces para nuestro proyecto, la enorme base de conocimientos y la diversidad de herramientas creadas en el universo Python decanta la balanza hacia ese lado. Quizás JavaScript permita mejores resultados usando la gestión de memoria en \textit{requests} simultáneas, pero a costa de un código más oscuro y difícil de mantener. Aunque JavaScript cuenta con un gran repertorio de paquetes Node.JS como utilidades de \textit{web-scraping}, en el entorno Python es difícil imaginar una tarea para la que no se haya escrito una (o más) herramientas que resuelvan eficazmente nuestro problema. 

Por otro lado, la comunidad de programadores de Python es inmensa y su creciente popularidad facilita el hallazgo de soluciones rápidamente, tanto en los foros como en la extensa documentación con la que cuenta. Aunque encontramos un rendimiento ligeramente inferior a otros lenguajes en ciertas búsquedas, es el precio a pagar por el tipado dinámico. 

Como valor añadido, Python es fácil de mantener cuando necesitamos adaptar nuestro código a las cambiantes estructuras de las páginas web. Además, sus reconocidas herramientas de análisis de datos nos permiten continuar en el mismo entorno, sin necesidad de buscar alternativas para afrontar tareas relacionadas con la \textit{data science}.

\section{Bibliotecas}
A lo largo del proyecto se ha recurrido a diversas bibliotecas de Python. A continuación, se presenta brevemente cada una de ellas.

\subsection{Scholarly}
Se trata de una biblioteca de Python que permite acceder a los datos de Google Scholar de manera fácil y rápida. La biblioteca proporciona una interfaz sencilla para buscar y recuperar información sobre artículos, autores y revistas en Google Scholar, incluyendo metadatos, citas y otra información relacionada.

Sin embargo, esta biblioteca ha terminado siendo descartada para este proyecto. A diferencia de otras técnicas de \textit{web-scraping}, Scholarly no está diseñada para extraer grandes cantidades de datos de Google Scholar. La biblioteca tiene una serie de limitaciones en cuanto a la cantidad de datos que se pueden recolectar, ya que está diseñada para ser utilizada en investigaciones científicas y no para la extracción masiva de datos. Además, Scholarly está diseñada para respetar los términos de servicio de Google Scholar y no violar la política de uso de la plataforma, por lo que no se recomienda su uso para recolectar grandes cantidades de datos.

En resumen Scholarly es una herramienta útil para acceder a información científica y académica de manera rápida y sencilla, pero no esta diseñada para extraer grandes cantidades de datos.

\subsection{Beautiful Soup}
Beautiful Soup es una biblioteca de Python extremadamente útil para la extracción de datos de páginas HTML o XML. Actualmente se encuentra en su versión 4.8.1. En el siguiente \href{https://beautiful-soup-4.readthedocs.io/en/latest/}{enlace} se puede acceder a la documentación de su página oficial. 

Esta biblioteca nos proporciona numerosos módulos para navegar a través de las páginas web y para extraer fácilmente su contenido. Puesto que gran parte del proyecto se basa en el uso de \textit{web-scraping}, se ha hecho uso intensivo de la misma para extraer los principales datos de los artículos científicos que posteriormente conformarán la BBDD del proyecto.

De Beautiful Soup hay que destacar su facilidad de uso y la amplia documentación que aporta. En su página web nos aconseja utilizar el analizador \textit{lxml}, que proporciona al entorno Python la disponibilidad de las bibliotecas \textit{libxml2} y \textit{libxslt}. Allí mismo, se anima incluso a utilizar aisladamente este parser cuando el tiempo de respuesta sea una cuestión crítica. En nuestro caso, las facilidades que proporciona Beautiful Soup justifican ampliamente su uso, aunque la  rapidez de resultados no iguale la de la utilización aislada de los analizadores sobre los que trabaja.
 

\subsection{Habanero}
Para poder hacer uso de la \textbf{API de Crossref}, se han explorado diversas alternativas. De entre todas las bibliotecas de Python se ha seleccionado Habanero, ya que es una biblioteca muy fácil de usar y está en constante actualización.

Esta biblioteca está diseñada para facilitar el acceso a las bases de datos de revistas científicas y a otras relacionadas con el ámbito académico. Ofrece una interfaz simple para recuperar información utilizando los protocolos y las API de diferentes bases de datos, incluyendo JSTOR, Unpaywall, Crossref, DataCite, etc. Además, Habanero es compatible con las normas de Open Access, lo que permite a los usuarios acceder a contenido científico gratuito y de libre acceso.

\subsection{Selenium}

Selenium es una biblioteca de Python de automatización de navegadores web que se utiliza para controlar el comportamiento de un navegador en tiempo real. Es una herramienta muy poderosa para realizar \textit{web scraping} de sitios web, especialmente en aquellos casos en que se necesita interactuar con la página web como si lo hiciera un usuario humano.

A diferencia de otras bibliotecas de \textit{web scraping}, Selenium permite simular la navegación humana en la web, permitiendo realizar búsquedas, hacer clic en botones, llenar formularios, entre otras acciones. Esto es muy útil cuando se trata de sitios web que tienen medidas de seguridad para evitar el \textit{web scraping} y requieren de interacción humana.

\subsection{Scikit-learn}

Scikit-learn es una librería de aprendizaje automático en Python que ofrece una amplia variedad de algoritmos de clasificación, regresión y clustering, entre otros. Esta librería destaca por su facilidad de uso y la implementación eficiente de los algoritmos, lo que permite su uso en grandes conjuntos de datos. Cuenta también con herramientas de preprocesamiento de datos y evaluación de modelos. Sus ventajas incluyen una documentación completa y una comunidad activa que ofrece soporte y mejora continuamente la librería, así como la capacidad de integrarse fácilmente con otras librerías y herramientas de Python para el análisis de datos.

\subsection{Unittest}

Para los proyectos de pruebas, hemos utilizado la biblioteca de pruebas unitarias de Python llamada unittest. Una de las principales ventajas de unittest es que se integra fácilmente con el flujo de trabajo de Test Data Driven (metodología de pruebas empleada), permitiendo que las pruebas se ejecuten automáticamente y de manera aislada del código principal. Además, unittest proporciona una gran cantidad de funcionalidades y herramientas para facilitar la escritura y ejecución de pruebas, lo que nos ha permitido realizar pruebas complejas de manera eficiente y eficaz.

\section{Bases de datos}

Antes de su compra por Oracle Corporation (2010), MySQL era la aplicación de base de datos más popular de código abierto para la programación web. En el momento actual, el software libre nos ofrece dos soluciones ampliamente contrastadas en gestores de bases de datos relacionales: MariaDB y PostgreSQL. Necesitaremos una herramienta de este tipo para organizar los datos recolectados. La estructura tabular de la información nos permite aplicar sobre ella el lenguaje de interrogación SQL.

En realidad, MariaDB es una bifurcación de MySQL nacida para garantizar la supervivencia del proyecto como código abierto. Hoy en día, MariaDB es altamente compatible con MySQL e incluso superior en sus últimas versiones (10.1.1), ya que la comunidad ha ido añadiendo nuevas características al proyecto original. Por otro lado, aunque no está tan extendido como MySQL, PostgreSQL es posiblemente el gestor de bases de datos de código abierto más sólido y potente a día de hoy.


\tablaSmall{Comparativa entre MariaDB y PostgreSQL}{ll}{comparativabbdd}
{MariaDB & PostgreSQL\\}{ 
No totalmente compatible con SQL. & Compatible con SQL estándar.\\
Soporte para tipos de datos estándar SQL. & Soporta además tipos avanzados.\\
Tipos de datos flexibles. & Tipos de datos estrictos. \\
Tamaño pequeño de base de datos. & Tamaño grande de base de datos. \\
Sin soporte directo para JSON.  & Soporte directo de JSON. \\
No índices parciales. & Índices parciales. \\
No soporte para web dinámica. & Soporta sitios web dinámicos.\\
} 


\subsection{PostgreSQL}
Finalmente se opta por \href{https://www.postgresql.org/}{PostgreSQL} (también denominado Postgres). Pese a que ambas opciones son muy similares, se ha elegido esta última debido a que ya se ha trabajado con ella anteriormente. Además, PostgreSQL posee una sólida reputación por su arquitectura comprobada, confiabilidad, integridad de datos, conjunto sólido de características, extensibilidad y la dedicación de la comunidad de código abierto detrás del software para ofrecer soluciones innovadoras y de rendimiento constante~\cite{Postgresql.org}.

Al igual que MariaDB, es un sistema de gestión de bases de datos relacional de código abierto. Así pues, está dirigido y desarrollado por una comunidad altruista de desarrolladores (PostgreSQL Global Development Group).
PostgreSQL utiliza y amplía el lenguaje SQL combinado con muchas características que almacenan y escalan de forma segura las cargas de trabajo de datos más complicadas~\cite{Postgresql.org}. 
Su última versión y la usada para el proyecto es la versión 15.1 lanzada el 10 de noviembre de este año.


\section{APIs}
A lo largo del proyecto se ha recurrido a varias APIs diferentes, todas ellas de acceso gratuito.

\subsection{Google Scholar}
Se trata de una de las principales herramientas que ofrece Google a los investigadores. \href{https://scholar.google.com/}{Google Scholar} es, fundamentalmente, un buscador de contenido y bibliografía científica que permite localizar artículos de revistas especializadas ordenados por relevancia en función de las palabras clave introducidas en el buscador. También se puede  filtrar la información en función de su fecha de publicación, idioma o número de citas.


\subsection{Crossref}

 Crossref es una organización sin ánimo de lucro que provee de una \href{https://www.crossref.org/}{herramienta} que facilita el acceso a la información de los artículos científicos a partir de su DOI\footnote{El DOI (\textit{Digital Object Identifier}) es el acrónimo con el que se conoce al identificador inequívoco y persistente de un objeto digital (como un artículo de revista, un libro electrónico, una imagen, etc...). Se trata de un enlace permanente al contenido electrónico de dicho objeto. Por lo general, un DOI tiene forma de código alfanumérico.}.

Como ya se ha mencionado previamente, en base a la información de los DOIs la agencia de registro CrossRef es capaz de proporcionarnos aplicaciones útiles para hacer nuestro flujo de investigación más sencillo.
Se trata de una asociación sin ánimo de lucro de editoriales científicas que no solo facilita el registro de DOIs a las editoriales, sino que también ofrece servicios y aplicaciones para el personal investigador que tienen como base estos códigos.

Los DOIs que CrossRef almacena van acompañados de información que refleja las cualidades básicas de una publicación científica. Me estoy refiriendo a datos como títulos, abstracts, palabras clave, autores…

CrossRef Metadata Search hace posible obtener toda esta información al instante con tan solo proporcionar el DOI asociado a una publicación, o al contrario, obtener el DOI de la publicación con tan solo introducir algunos de estos datos en su buscador.









\capitulo{5}{Aspectos relevantes del desarrollo del proyecto}

%Este apartado pretende recoger los aspectos más interesantes del desarrollo del proyecto, comentados por los autores del mismo.
%Debe incluir desde la exposición del ciclo de vida utilizado, hasta los detalles de mayor relevancia de las fases de análisis, diseño e implementación.
%Se busca que no sea una mera operación de copiar y pegar diagramas y extractos del código fuente, sino que realmente se justifiquen los caminos de solución que se han tomado, especialmente aquellos que no sean triviales.
%Puede ser el lugar más adecuado para documentar los aspectos más interesantes del diseño y de la implementación, con un mayor hincapié en aspectos tales como el tipo de arquitectura elegido, los índices de las tablas de la base de datos, normalización y desnormalización, distribución en ficheros3, reglas de negocio dentro de las bases de datos (EDVHV GH GDWRV DFWLYDV), aspectos de desarrollo relacionados con el WWW...
%Este apartado, debe convertirse en el resumen de la experiencia práctica del proyecto, y por sí mismo justifica que la memoria se convierta en un documento útil, fuente de referencia para los autores, los tutores y futuros alumnos.

En esta sección se presenta un resumen de los hallazgos más relevantes obtenidos a lo largo del proyecto, así como una descripción detallada de las distintas fases por las que se ha atravesado hasta lograr alcanzar una solución satisfactoria al problema planteado.

\section{Extracción de la información}
La primera fase de este proyecto es la más extensa y consiste en la extracción de datos para alimentar la base de datos. Para esto, es necesario tener en cuenta varias bases de datos bibliográficas que ofrecen diversos recursos para los investigadores. Entre las más importantes se encuentran Google Scholar (GS), Scopus, Web of Science (WoS) y PubMed.

\subsection{Prototipos para Google Scholar}

De todas las bases de datos bibliográficas mencionadas, la más destacada es Google Scholar, puesto que crece más rápido que cualquier base de datos tradicional en todos los campos científicos~\cite{harzing2010}~\cite{lopez2017}.

Así pues, el primer prototipo diseñado consiste en la extracción de datos de GS. Sin embargo, si bien es cierto que GS es un motor de búsqueda gratuito, universal y rápido con una amplia cobertura, tiene muchas limitaciones~\cite{lopez2017}. 


Por ejemplo, no hay funcionalidades de exportación de datos o API disponibles debido a restricciones comerciales. Además, solo se pueden mostrar los primeros 1,000 resultados de cada consulta. Aunque algunas técnicas como los retrasos temporales o el uso de \textit{proxies} pueden ayudar, no resuelven completamente estas limitaciones. Además, la extracción automatizada de estos datos va en contra de las políticas del archivo robots.txt de GS, lo que hace que los usuarios que realizan demasiadas consultas automatizadas sean bloqueados cada 200 solicitudes detectadas~\cite{lopez2017}.

Durante el avance del proyecto, nos dimos cuenta de muchas de las limitaciones mencionadas.

\subsubsection{Prototipo inicial}

Tras comprender que las complicaciones eran numerosas, se decidió crear un \textbf{prototipo sencillo}, consistente en un \textit{script} en Python, que trataría de lanzar mil peticiones de búsqueda. Esto nos permitiría establecer los límites de realizar \textit{web-scrapping} sobre Google Scholar.

Así pues, manos a la obra, se desarrolló un \textit{script} sencillo, que solicita acceso a la página principal de Google Scholar y, mediante métodos HTTP, realiza una búsqueda (a partir de parámetros solicitados por pantalla). Finalmente, extrae el título de la página resultante tras hacer la búsqueda. 

Todo esto se logra haciendo uso de la biblioteca de Python BeautifulSoup (ver documentación en el siguiente \href{https://beautiful-soup-4.readthedocs.io/en/latest/}{enlace}).

% Filtro por revistas

El primer obstáculo encontrado es que Google Scholar no permite hacer búsquedas en función de la revista. Si se trata de escribir el nombre de una revista en el buscador, aparecerán artículos relacionados que mencionan esa revista, pero no siempre artículos de la revista en cuestión. En cambio, permite hacer búsquedas a partir de palabras clave.
Sin embargo, aunque no se puede buscar directamente por revista, sí se puede extraer de los resultados encontrados tras buscar una temática concreta.

\imagen{ejemploResultadoGS}{Ejemplo de búsqueda por palabra clave}{1}

Sin embargo, existe una problemática: en cada artículo se hace referencia a la revista que lo publica de forma distinta. Por ejemplo, como se puede apreciar en la siguiente imagen, la revista se encuentra ubicada en una <<etiqueta>> HTML diferente.

\imagen{etiquetaDeRevista}{Localización de la revista en Google Scholar}{1}

Por ello se ha buscado una solución alternativa.
Existe una etiqueta (\texttt{<b> … </b>}) que incluye aquellas palabras resaltadas en negrita en la página web. Google Scholar resalta en una búsqueda aquellas palabras que coinciden con alguno de los parámetros buscados. Así pues, si se realiza una búsqueda en Google Scholar pasando como parámetro el nombre de la revista que nos interesa, bastará con buscar en la clase  \texttt{gs\_a} (donde se almacenan los detalles del artículo) alguna etiqueta \texttt{<b>}.
Se ha escrito un pequeño código de prueba y, en principio, parece que funciona sin problemas. Se adjunta dicho código (extraerRevistas.py) en el repositorio de GitHub.
Este código hace dos búsquedas: una a la revista Alas Peruanas y otra a la revista The Lancet (elegidas de forma aleatoria).

%La salida obtenida (solo de la primera página de resultados) es la siguiente:

%\imagen{salidaObtenida}{Salida por consola}{1}

%La salida por pantalla mostrada en la imagen anterior sigue la siguiente estructura: [\texttt{‘nombre artículo’, ‘nombre de la revista’, número\_citas, fecha, id}]
La salida que se mostrará por pantalla sigue la siguiente estructura: [\texttt{‘nombre artículo’, ‘nombre de la revista’, número\_citas, fecha, id}]

De igual forma se recogen los resultados en los CSVs (codificados en UTF-8) con nombre: BBDD1.csv y BBDD2.csv.


% Paginación

Otra de las mejoras que se implementa en el prototipo en la capacidad de navegar a través de las distintas páginas de resultados de Google Scholar. Para ello, se incluye un nuevo campo en la \textit{query} de búsqueda que irá aumentando de 10 en 10 por cada página de resultados nueva que se consulte.
Ahora ya se obtiene un número de resultados considerable (aproximadamente 100). Estos resultados se guardarán también en los CSVs (también codificados en UTF-8) con los siguientes nom-bres: BBDD3.csv y BBDD4.csv

% CAPTCHA y proxies

Tras diseñar y programar el código mencionado, se procede a su prueba. La primera ejecución del \textit{script} resultó desaletadora, ya que, a partir de la solicitud número \textbf{726}, Google Scholar detecta que un \textbf{\textit{bot}} está realizando búsquedas. A partir de ese momento, nuestra dirección IP queda bloqueada y las solicitudes fallan sin excepción. Google Schoolar nos redirecciona a una página (figura \ref{fig:error}) donde se solicita al usuario resolver un \textit{captcha}. 

\imagen{error}{Captura de reCAPTCHA de Google}{1}

El número de solicitudes exitosas es demasiado bajo para cumplir su función en nuestro proyecto, por lo que se procede a buscar una solución alternativa. 

Tras distintas pruebas, se encontró una forma de superar la barrera del \textit{captcha}. A saber: añadiendo a la \textit{url} una sección de texto extra que permite suprimir esta excepción durante un periodo concreto de tiempo. Así pues, logramos ejecutar con éxito el \textit{script} tantas veces como fuese necesario. Sin embargo, esta solución tampoco es válida a largo plazo. Si se intenta ejecutar de nuevo el programa, en otra sesión, la dirección vuelve a ser inválida. Además, es un remedio poco práctico, ya que se debe concatenar distintos parámetros y cadenas de texto que, dependiendo del momento y el contexto en que se ejecute, pueden no servir.

Como la propuesta anterior no fue satisfactoria, se siguió buscando opciones. La siguiente propuesta consistía en usar un agente de usuario distinto para cada solicitud\footnote{Un agente de usuario es cualquier software, que actúa en nombre de un usuario, que <<recupera, presenta y facilita la interacción del usuario final con el contenido web>>. Algunos ejemplos destacados de agentes de usuario son los navegadores web. La cadena User-Agent es uno de los criterios por los cuales los rastreadores web pueden ser excluidos del acceso a ciertas partes de un sitio web utilizando el Estándar de exclusión de robots (\texttt{archivo robots.txt}).}. De esta forma <<enmascaramos>> nuestra dirección IP. La biblioteca \textbf{Request} de Python ofrece métodos para lograrlo. Dicho esto, se implementó un método que extrae \textit{proxies} de listas públicas y gratuitas de Internet (e.g.: \href{https://www.proxyscrape.com/free-proxy-list}{proxyscrape.com}). Se prueba la ejecución del prototipo una vez más y, finalmente, funciona sin inconvenientes. Ahora ya se puede decir que el proyecto es \textbf{viable}.


% DOI y Crossref
\subsubsection{Prototipo con extracción del DOI}

La siguiente meta de nuestro prototipo es la obtención del DOI, que es un campo fundamental que funcionará a modo de clave primaria en la base de datos. De esta forma se busca evitar futuros errores a la hora de identificar un artículo o a la hora de comparar artículos para eliminar duplicados. Sin embargo, se plantea una problemática al respecto: Google Scholar no comparte el DOI de los artículos en ninguna división de su página web (si bien es cierto que algunos artículos lo incluyen al final de su URL, pero no siempre ocurre esto). Puesto que por el momento no existe una forma de obtenerlo directamente, se han propuesto varias soluciones, a saber:
\begin{itemize}
	\item Comprobar el porcentaje de éxito de extraer el DOI de la URL de los artículos (en aquellos casos en los que aparezca).
	\item Acceder a la página del artículo en cuestión y extraerlo de dicha página.
         \item Introducir el nombre del artículo en la página de \href{https://www.crossref.org/}{Crossref} o similares, que te permiten obtener el DOI del artículo cuyo nombre pases como parámetro.
\end{itemize}
Como conclusión, se descarta la primera opción tras hacer una breve prueba, ya que falla en seguida. Se descarta también la segunda opción, ya que cada página sitúa el DOI en un sitio haciendo muy difícil la búsqueda del mismo a través de un algoritmo. Por lo tanto, la opción más adecuada en este caso es la tercera. Aunque supone una búsqueda extra en la complejidad algorítmica del prototipo, es la única forma segura de calcular el DOI sin equivocaciones. Por lo tanto, se incluye en el bucle del prototipo un nivel extra de profundidad de búsqueda.

Para hacer la búsqueda en Crossref es necesario indicarle tanto el título como el año de publicación para asegurarnos de que se obtiene exactamente el artículo que deseamos y no otros similares. Además, es preciso tener en cuenta que Crossref no permite a los programas hacer búsquedas desde la misma interfaz que los usuarios (para evitar que los programas bloqueen las búsquedas de los usuarios). Por lo tanto, se realizarán las búsquedas en la dirección donde se ubica la API creada por Crossref especialmente para la extracción de datos por parte de programas.

% Multihilo
\subsubsection{Prototipo multihilo}

Con el prototipo listo, se prueba a extraer la información de los artículos de las primeras 20 páginas de resultados de Google Scholar de dos revistas. Los resultados no son muy esperanzadores: se obtienen 180 artículos de cada revista en 15 minutos. Estos resultados no son eficientes, por lo que se tratará de optimizar el prototipo siguiendo dos pautas. Primero, se tratará de extraer las llamadas a Crossref de forma que solo se tenga que realizar una única llamada una vez que se han extraído el resto de detalles sobre los artículos. La segunda pauta es emplear programación concurrente para ejecutar varios hilos al mismo tiempo. 

Así, cada revista será procesada por un hilo distinto. Para ello, será necesario recurrir a la librería de Python \textit{multiprocessing}.Tras actualizar estos cambios, el prototipo obtiene resultados mucho más óptimos: obtiene los 180 artículos de ambas revistas en tan solo 8 minutos. El próximo objetivo es determinar el número máximo de hilos que se pueden lanzar sin impactar negativamente el rendimiento del prototipo.

Para ello, se solicitó permiso para acceder a los servidores de la universidad. A través de SSH, nos conectamos a una de las máquinas y, utilizando un entorno virtualizado de Miniconda, comenzamos a lanzar hilos. Sin embargo, el resultado fue desastroso, ya que al realizar 200 llamadas a GS, nos bloqueaban y teníamos que resolver un CAPTCHA. Intentamos insertar esperas de tiempo aleatorias y cambiar el user-agent en cada una de las llamadas. Incluso se probó a cambiar el término de búsqueda. Sin embargo, nada de esto funcionó. Matemáticamente, después de aproximadamente 200-250 llamadas, aparecía el CAPTCHA.


\subsubsection{Prototipo con Selenium}
La única solución posible frente a las restricciones de GS es tratar de <<humanizar>> las búsquedas para que parezca que las realiza un ser humano. Por lo tanto, se decidió emplear Selenium, que nos permite emular los pasos que realizaría un usuario normal al hacer una búsqueda. El nuevo prototipo entra directamente en el navegador con la ruta de la búsqueda, lo que ahorra la necesidad de realizar llamadas adicionales. Luego, hace clic en la sección citar de cada uno de los resultados de la búsqueda y aparece un cuadro emergente del cual se puede extraer la información completa de la cita. Esto resuelve otro de los problemas que teníamos con el prototipo anterior: GS acorta los nombres largos añadiendo puntos suspensivos, mientras que en la sección citar aparecen los nombres completos. 

Después de extraer esta información, se cierra el cuadro y se procede a realizar la misma acción en el siguiente artículo. Cuando se llega al décimo resultado (ya que solo hay 10 artículos por página de resultados), se emula un clic en el botón <<Siguiente>> y se procede a realizar lo mismo en la siguiente página de resultados, todo con pausas aleatorias de tiempo entre cada acción. Posteriormente, se procesan los datos de las citas usando expresiones regulares y transformando los diccionarios de datos en un CSV.

Como nota adicional, cabe mencionar que, si se incluye la etiqueta \texttt{source: <nombre de la revista>} en la búsqueda, es posible hacer una búsqueda solo por el nombre de la revista, algo que desconocíamos hasta ahora.

Aunque este método ha probado ser mejor que el anterior, ya que se obtuvieron casi 300 resultados, GS detecta que son consultas automatizadas y salta el CAPTCHA al final.

Para intentar resolver esta situación, se ha intentado superar el CAPTCHA emulando un clic sobre el botón correspondiente. No obstante, para lograrlo es necesario superar varias capas ocultas que se encuentran sobre dicho botón, lo cual resulta complejo. Una vez se logra acceder al botón, aparecen las imágenes que se deben reconocer y esto no se puede automatizar de manera sencilla. Se llega a la conclusión de que no es posible forzar más llamadas de las permitidas en Google Scholar, tal y como se ha mencionado en otros trabajos, como en el caso de Publish or Perish~\cite{harzing2010}.



\subsection{Prototipos para Crossref}
% TODO

Vista la imposibilidad de obtener una cantidad aceptable de resultados de Google Scholar, se decidió cambiar la fuente de datos. En este caso, Crossref, aunque no es una base de datos tan potente, tiene una API gratuíta y sin tantas restricciones.



\capitulo{6}{Trabajos relacionados}

Este apartado sería parecido a un estado del arte de una tesis o tesina. En un trabajo final grado no parece obligada su presencia, aunque se puede dejar a juicio del tutor el incluir un pequeño resumen comentado de los trabajos y proyectos ya realizados en el campo del proyecto en curso. 

\capitulo{7}{Conclusiones y Líneas de trabajo futuras}

%Todo proyecto debe incluir las conclusiones que se derivan de su desarrollo. Éstas pueden ser de diferente índole, dependiendo de la tipología del proyecto, pero normalmente van a estar presentes un conjunto de conclusiones relacionadas con los resultados del proyecto y un conjunto de conclusiones técnicas. 
%Además, resulta muy útil realizar un informe crítico indicando cómo se puede mejorar el proyecto, o cómo se puede continuar trabajando en la línea del proyecto realizado. 


\section{Líneas de trabajo futuras}
En un futuro, si se desea utilizar la aplicación para calcular el índice de impacto en un campo distinto al de \textit{computer science}, solo se requiere subministrar (con permisos de administrador de la API) un CSV con la lista de revistas del nuevo campo deseado. La red se volverá a entrenar y se reiniciará todo el proceso, permitiendo así la adaptación de la aplicación a diferentes áreas de investigación. Con esta flexibilidad, la aplicación puede ser empleada para el cálculo del índice de impacto en una amplia variedad de campos, ampliando así su alcance y aplicabilidad.
\capitulo{8}{Prototipo inicial}

Al revisar la viabilidad del proyecto se plantearon posibles dificultades que podían surgir. Tras comprender que las complicaciones eran numerosas, se decidió crear un \textbf{prototipo sencillo}, consistente en un \textit{script} en Python, que trataría de lanzar  mil peticiones de búsqueda. Esto nos permitiría establecer los límites de realizar ``web scrapping'' sobre Google Scholar .

Así pues, manos a la obra, se desarrolló un \textit{script} sencillo, que solicita acceso a la página principal de Google Scholar y, mediante métodos HTTP, realiza una búsqueda (a partir de parámetros solicitados por pantalla). Finalmente, extrae el título de la página resultante tras hacer la búsqueda. 

Todo esto se logra haciendo uso de la biblioteca de Python \textbf{BeautifulSoup} (ver documentación en el siguiente \href{https://beautiful-soup-4.readthedocs.io/en/latest/}{enlace}). Esta biblioteca contiene métodos centrados en la extracción de datos de archivos HTML y XML para su posterior análisis. Es fácil advertir cuán idónea resulta esta biblioteca para nuestro propósito.

Tras diseñar y programar el código mencionado, se procede a su prueba. La primera ejecución del \textit{script} resultó desaletadora, ya que, a partir de la solicitud número \textbf{726}, Google Scholar detecta que un \textbf{\textit{bot}} está realizando búsquedas. A partir de ese momento, nuestra ip queda bloqueada y las solicitudes fallan sin excepción. Google Schoolar nos redirecciona a una página (figura \ref{fig:error}) donde se solicita al usuario resolver un \textit{captcha}. 

\imagen{error}{Captura de reCAPTCHA de Google}{.70}

El número de solicitudes exitosas es demasiado bajo para cumplir su función en nuestro proyecto, por lo que se procede a buscar una solución alternativa. 

Tras distintas pruebas, se encontró una forma de superar la barrera del \textit{captcha}. A saber: añadiendo a la \textit{url} una sección de text extra que permite suprimir esta excepción durante un periodo concreto de tiempo. Así pues, logramos ejecutar con éxito el \textit{script} tantas veces como fuese necesario. Sin embargo, esta solución tampoco es válida a largo plazo. Si se intenta ejecutar de nuevo el programa, en otra sesión, la dirección vuelve a ser inválida. Además, es un remedio poco práctico, ya que se debe concatenar distintos parámetros y cadenas de texto que, dependiendo del momento y el contexto en que se ejecute, pueden no servir.

Como la propuesta anterior no fue satisfactoria, se siguió buscando opciones. La siguiente propuesta consistía en usar un agente de usuario distinto para cada solicitud

agente de usuario es cualquier software, que actúa en nombre de un usuario , que "recupera, presenta y facilita la interacción del usuario final con el contenido web". Por lo tanto, un agente de usuario es un tipo especial de agente de software .
Algunos ejemplos destacados de agentes de usuario son los navegadores web
La cadena User-Agent es uno de los criterios por los cuales los rastreadores web pueden ser excluidos del acceso a ciertas partes de un sitio web utilizando el Estándar de exclusión de robots ( archivo robots.txt ).

usar un \textit{proxy} para ``enmascarar'' nuestra ip. La biblioteca \textbf{Request} de Python ofrece métodos para lograrlo. Dicho esto, se implementó un método que extrae \textit{proxies} de listas públicas y gratuitas de Internet (v.g.: \href{https://www.proxyscrape.com/free-proxy-list}{proxyscrape.com}). Se prueba la ejecución del prototipo una vez más y, finalmente, funciona sin inconvenientes. Ahora ya se puede decir que el proyecto es \textbf{viable}.






...


- Incluir doi de crossref y resultados (obteniendo unos 170 artículos en 15 minutos)

Estos resultados no son eficientes, por lo que se tratará de optimizar el prototipo siguiendo dos pautas. Primero, se tratará de extraer las llamadas a Crossref de forma que solo se tenga que realizar una única llamada una vez que se han extraído el resto de detalles sobre los artículos. La segunda pauta es emplear programación concurrente para ejecutar varios hilos al mismo tiempo. 

Para poder realizar ambas cosas, se separarán las peticiones a Google Schoolar y las peticiones a Crossref en dos hilos distintos. Para ello, será necesario recurrir a la librería de Python \textbf{multiprocessing}.

Tras actualizar estos cambios, el prototipo obtiene resultados mucho más óptimos: artículos en 15 minutos.





\bibliographystyle{plain}
\bibliography{bibliografia}

\end{document}
